\documentclass[conference]{IEEEtran}
\IEEEoverridecommandlockouts

\usepackage{cite}
\usepackage{amsmath,amssymb,amsfonts}
\usepackage{graphicx}
\usepackage{textcomp}
\usepackage{xcolor}
\usepackage[utf8]{inputenc}
\usepackage[T1]{fontenc}

\def\BibTeX{{\rm B\kern-.05em{\sc i\kern-.025em b}\kern-.08em
    T\kern-.1667em\lower.7ex\hbox{E}\kern-.125emX}}

\begin{document}

\title{lorem ipsum}

\author{
  \IEEEauthorblockN{Habb}
  \IEEEauthorblockA{
    \textit{Department of ...}\\
    University\\
    City, Country\\
    email@domain.com
  }
}

\maketitle

\begin{abstract}
TODO: abstract.
\end{abstract}

\begin{IEEEkeywords}
keyword1, keyword2, keyword3
\end{IEEEkeywords}

\section{Introduction}
% \cite{TESTKEY}

\section{Related Work}

\section{Methodology}

\section{Experiments and Results}

\section{Conclusion}

% \nocite{*} menampilkan SEMUA entri refs.bib supaya template ini langsung
% ter-compile walau belum ada \cite. Hapus baris ini begitu mulai menulis dan
% mengutip referensi sungguhan dengan \cite{citekey}.
\nocite{*}
\bibliographystyle{IEEEtran}
\bibliography{refs}

\end{document}
